\documentclass{article}
\title{ECE 496 -- Homework 3}
\usepackage{graphicx}
\usepackage{pdfpages}
\usepackage{fancyhdr}
\usepackage{enumitem}
\usepackage{pdfpages}
\usepackage{amsmath}
\usepackage{amssymb}
\usepackage{float}
\usepackage[margin=1in]{geometry}
\pagestyle{fancy}
\fancyhf{}
\fancyhead[L]{Gianni Avilan-Losee}
\fancyhead[C]{ECE 496}
\fancyhead[R]{Homework 3}
\fancyfoot[C]{\thepage}
\usepackage{microtype}
\usepackage{textcomp, gensymb}
\author{Gianni Avilan-Losee}
\begin{document}
\setcounter{section}{2}
\setcounter{subsection}{1}
\begin{enumerate}[label=\thesubsection]
	\item
	\begin{enumerate}[label=\alph*)]
	\item At each mask step, there is a fraction of the die that are good. At the next mask step, a fraction of the remaining die will be good, and so on. Since the potential for errors compounds, to get a total fraction of 0.3 good die at the end, each mask step should produce, at minimum, a fraction of \(0.3^{\frac{1}{25}}=0.953\) good die.
	\item Similarly, for a total fraction of 0.7 good die, each mask step should produce a minimum fraction of \(0.7^{\frac{1}{25}} = 0.9858\) good die.
\end{enumerate}
\setcounter{subsection}{4}
\item
	\begin{enumerate}[label=\alph*)]
		\item For the n-diffusion window, negative photoresist should be used to create a window in the photoresist layer, where the exposed area of barrier layer can be easily etched away with a solvent, and then n-doped. Then, the contact windows can also be etched away using a negative photoresist. Finally, a positive photoresist can be used for the metal etch step, which \textit{leaves} a layer of barrier material in the exposed area. Then, the metal etch can be applied, and will only contact the exposed substrate material where there is no barrier layer.
		\item The textbook contains an example of potential alignment marks for a multi-mask process, shown below. It's likely that a similar pattern could be used in the above process.
			\begin{figure}[H]
			  \centering
			  \includegraphics[width=0.4\textwidth]{homework-3_fig-1.jpg}
			  \caption{Potential alignment marks for the diode process in Q-2.4(a)}
			  \label{fig:alignment}
			\end{figure}
	\end{enumerate}
	\setcounter{subsection}{6}
	\item
		\begin{enumerate}[label=\alph*)]
			\item For a feature size of \(0.25 \ \micro\mathrm{m}\) with a numerical aperture of \(1\), using 
				\[
				  F = 0.5 \frac{\lambda}{NA},
				\]
				\(\lambda = \frac{F \times NA}{0.5} = \frac{0.25 \times 10^{-6}}{0.5} = 250 \ \mathrm{nm}\).

				Then, depth of field can be calculated using 
				\[
				  DF = 0.6 \frac{\lambda}{(NA)^2},
				\]
				as \(DF = 0.6 \times 250 \times 10^{-9} = 150 \ \mathrm{nm}\).
			\item Repeating the entire process for \(NA = 0.5\), \(\lambda = \frac{0.25 \times 10^{-6} \times 0.5}{0.5} = 125 \ \mathrm{nm}\) and \(DF = \frac{0.6 \times 125 \times 10^{-9}}{0.25} = 300 \ \mathrm{nm}.\)
		\end{enumerate}
		\setcounter{subsection}{7}
	\item Using \(NA = 1\), the largest aperture size, and solving for feature size, \(F = 0.5 \frac{193 \times 10^{-9}}{1} = 96.5 \ \mathrm{nm}\).
		\setcounter{subsection}{8}
	\item Using the process for the previous \(F = 0.5\frac{13 \times 10^{-9}}{1} = 6.5 \ \mathrm{nm}\).
\end{enumerate}
\end{document}
